\section{Products of curves and surfaces and their Chow summands}\label{sec:products}

\subsection{The product-model input}

The product theorem uses a numerical calculation on a particular
semistable model. We state its algebraic content before applying the
companion comparison with actual-form obstructions.

Let $R$ be a complete discrete valuation ring with finite residue field.
Take regular projective reduced strictly semistable models of finitely
many curves and surfaces. Fix orders on their components and form the
ordered staircase resolution of their product. Write its special fiber
as $D=\bigcup_iD_i$, with $\dim D_i=d$, retaining each connected
intersection separately. For an oriented edge $i\to j$, let
$j_i:E_{ij}\to D_i$ and $j_j:E_{ij}\to D_j$ be the inclusions. Set
\[
 (\rho x)_{ij}=j_j^*x_j-j_i^*x_i,\qquad
 \sigma y=(j_{i*}y,-j_{j*}y),\qquad A=\sigma\rho.
\]
These are cohomological actions of actual rational Chow correspondences.
The off-diagonal entries of $A$ have codimension $d+1$; its diagonal
entry on $D_i$ is
\[
 (\Delta_i)_*c_1(N_{D_i/\mathcal D})
       =-(\Delta_i)_*\sum_{j\ne i}[E_{ij}].
\]

\begin{proposition}[Companion product-model theorem]\label{prop:product-model}
For this ordered model there is an invertible prime $\ell$ and a
matrix of rational Chow correspondences $G$, with entries of
codimension $d-1$, for which
\begin{equation}\label{eq:inner-inverse}
 A^H G^H A^H=A^H
\end{equation}
on full geometric $\ell$-adic cohomology, in every degree and with
its Tate twist. This remains true after product with a smooth
projective $R$-scheme, by replacing $A,G$ with the exterior products
of their correspondences and the coefficient's diagonal.
\end{proposition}

The proof is \cite[model splitting theorem]{InternalProducts}, whose
actual deletion and residual-core constructions are included in the
proof dossier. The essential structure is as follows. Each component
is a tower of blowups of closed strata in the full reduced boundary
of a product of smooth curve and surface components. For a subset
$I$ of the surface indices, coordinate deletion gives an actual
proper birational map
\begin{equation}\label{eq:deletion-map}
 f_I:D_{\boldsymbol i}\longrightarrow
 \left(\prod_{j\in I}X_{j,i_j}\right)\times W_{-I,\boldsymbol i_{-I}},
\end{equation}
where $W_{-I}$ is the staircase component of the remaining factors.
The refined surface projectors $p_j$ onto $H^2_{\rm tr}$ annihilate
cohomology supported on the boundary curves. Consequently the transported
projectors
\[
 P_I=f_I^*((\boxtimes_{j\in I}p_j)\boxtimes1)f_{I*}
\]
satisfy $P_I=\prod_{j\in I}P_{\{j\}}$, commute with the incidence,
and do so on exceptional summands as well. The proof of commutation
uses support localization for the kernel of a birational pull-push
operator in \eqref{eq:deletion-map}.

The simultaneous sectors retain the tensor factors $t_2$ selected by
these projectors. The residual projector annihilates each generator
which still contains a full-surface $t_2$ factor; the deletion proof
establishes this equality in relative Chow before passing to
localization. The surviving blowup and localization generators
in the complementary factor are of abelian type. Thus the incidence
operator on a sector is the identity on its retained transcendental
factor tensored with an operator on one abelian carrier. Numerical
semisimplicity and the one-prime comparison for that carrier give its
inner inverse, which is tensored back with the retained factor.
Finite-field choices descend by normalized averaging: the equation
\eqref{eq:inner-inverse} is linear in $G$. This is the algebraic
product-model statement used here.

Two consequences of \eqref{eq:inner-inverse} explain its role in the
height problem. Let $\mathscr A^u$ be the space of actual algebraic
cycle classes on the disjoint union of the $D_i$, and let
$p:\mathscr A^u\twoheadrightarrow N^u(D^{(0)})_\Q\otimes\Q_\ell$
be the quotient to numerical classes. If $p(x)$ is in the numerical
kernel of $A$, then
\begin{equation}\label{eq:numerical-correction}
 x^\flat=x-G^HA^Hx
\end{equation}
has the same numerical image and belongs to $\ker A^H$. Also,
if an algebraic class $\alpha$ is in $\im A^H$, then $G^H\alpha$
is algebraic and $A^HG^H\alpha=\alpha$. Hence
\begin{equation}\label{eq:image-saturation}
 \mathscr A^{u+1}\cap\im A^H=A^H\mathscr A^u.
\end{equation}
The numerical-kernel lifting and actual-span saturation are the two
inputs to the companion model criterion \cite{InternalModel}.

\begin{theorem}[Arbitrary curve and surface products]\label{thm:products}
Let $X_1,\ldots,X_m/K$ be arbitrary smooth projective geometrically
integral curves or surfaces. Let $T/K$ be smooth projective
geometrically integral with potentially good projective reduction. Then
\[
 X=T\times\prod_{a=1}^mX_a
\]
has $\mathsf V$. The assertion holds over every finite extension of
$K$ for every newly specified geometrically homologically trivial
cycle there. Empty products and a point as coefficient are permitted.
\end{theorem}

\begin{proof}
After a common finite extension choose projective strictly semistable
alterations of the low-dimensional factors and a smooth projective
model of $T$. Semistable repair after any intermediate ramified
extension is part of the ordered-model construction in
\cite{InternalProducts}. Apply \propref{prop:product-model} to the
resulting witness. Weight--monodromy for curves and surfaces, the
purity of the good coefficient, and K\"unneth give the odd-weight
condition in the companion model criterion. Equations
\eqref{eq:numerical-correction} and \eqref{eq:image-saturation}
supply its two exactness conditions. The pulled-back obstruction
therefore vanishes. Dominant and finite-degree descent give zero
obstruction over $K$; the original curved pair gives the requested
form.

For an arbitrary finite $K'/K$, the factors still have the same
dimensions and the coefficient remains potentially good. Reapply
the construction over $K'$ to its newly specified cycle. This proves
the assertion for new classes, as well as for pullbacks of old ones.
\end{proof}

Over a number field, the curves and surfaces are unrestricted. If a
coefficient $T$ is present, it is required to have potentially good
projective reduction at every finite place. Applying
\thmref{thm:products} at all completions and using
\lemref{lem:two-degrees} gives the full height in every complementary
degree, also after every finite number-field extension.

\subsection{Abelian factors and the surface-generated Chow domain}

\begin{lemma}\label{lem:abelian-curve-source}
Over a characteristic-zero field, an abelian variety becomes a
rational Chow-motive retract of a product of curves after a finite
extension.
\end{lemma}

\begin{proof}
For an abelian variety of dimension $g\geq2$, choose a smooth
complete-intersection curve and a point $c_0$ on it after a finite
extension. Translate by $-c_0$ and call the resulting pointed curve
$C$; choose the complete intersection so that $C-C$ generates $A$.
The addition map $C^g\to A$ is dominant. Indeed,
the successive sums of $C$ increase in dimension until they fill
$A$: stabilization at a proper subvariety would make that subvariety
stable under the group generated by $C$. Since source and target
have dimension $g$, the map is generically finite, of some degree
$d>0$. Its graph satisfies $f_*f^*=d\Delta_A$, and
$d^{-1}f_*$ is a Chow retraction. Dimension one is immediate;
dimension zero gives the unit motive.
\end{proof}

\begin{corollary}\label{cor:surface-envelope}
Suppose that after a finite extension the rational Chow motive of
$X$ is a direct summand of a finite direct sum of Tate twists of
products of curves and surfaces. Then $\mathsf V(X)$ holds. It also
holds for $T\times X$ with a potentially good projective coefficient.
These conclusions include products with arbitrary abelian factors.
\end{corollary}

\begin{proof}
Choose the inclusion and retraction correspondences. Their composites
sum to the diagonal. Tensor with the coefficient's diagonal when
needed, apply \thmref{thm:products} to each intermediate variety,
and use \lemref{lem:general-transfer}. An abelian factor may first be
replaced by its curve-product retraction from
\lemref{lem:abelian-curve-source}.
\end{proof}

\subsection{Hilbert schemes of points on arbitrary surfaces}

Let $S$ be a smooth projective geometrically integral surface and
let $H=S^{[m]}$. For a
partition $\lambda=1^{a_1}\cdots m^{a_m}$, put
$l(\lambda)=\sum_i a_i$, $s_\lambda=m-l(\lambda)$, and
$G_\lambda=\prod_i\mathfrak S_{a_i}$. On the smooth power
$S^{l(\lambda)}$, use the rational projector
\[
 e_\lambda=|G_\lambda|^{-1}\sum_{g\in G_\lambda}\Gamma_g.
\]
The theorem of de Cataldo--Migliorini
\cite[Theorem 6.2.1 and Remark 6.2.2]{DCM} gives
\begin{equation}\label{eq:hilbert-motive}
 h(S^{[m]})\simeq\bigoplus_{\lambda\vdash m}
       (h(S^{l(\lambda)}),e_\lambda)\otimes\LL^{s_\lambda}.
\end{equation}
We use the covariant twist convention
$\CH_a(h(V)(s))=\CH_{a-s}(V)_\Q$.

\begin{theorem}\label{thm:hilbert}
For arbitrary smooth projective geometrically integral surfaces
$S_i$, curves $C_j$, and
positive integers $m_i$, the variety
\[
 T\times\prod_j C_j\times\prod_i S_i^{[m_i]}
\]
has $\mathsf V$ if $T$ is smooth projective geometrically integral
and has potentially good projective reduction.
Over a number field, it has a full height domain when this reduction
condition holds at every finite place.
\end{theorem}

\begin{proof}
Tensor the actual incidence correspondences in
\eqref{eq:hilbert-motive}. For a tuple of partitions they factor
through $T\times\prod_j C_j\times\prod_i S_i^{l(\lambda_i)}$.
If $s=\sum_i s_{\lambda_i}$, this intermediate variety has
dimension $\dim X-2s$ and an original codimension-$r$ class maps
to codimension $r-s$. Its homological class is zero. Apply
\thmref{thm:products} in every feasible degree and then
\lemref{lem:general-transfer}. The group in degree zero has no
homological kernel, and out-of-range groups vanish. The global
assertion follows at all completions. The correspondences commute
with base change, so the same proof applies to all new classes
after finite extension.
\end{proof}

\subsection{Finite specialization of Chow certificates}

\begin{lemma}\label{lem:specialize-certificate}
Let $X/k$ be smooth projective in characteristic zero. Suppose that
over a field extension $\Omega/k$ its diagonal factors in rational
Chow through finitely many products of smooth projective curves and
surfaces. Then such a factorization exists over a finite extension
$E/k$, with auxiliary curves and surfaces allowed to change.
\end{lemma}

\begin{proof}
Choose cycle representatives for the correspondences and finitely
many principal divisors witnessing the diagonal equality. All these
data, together with the auxiliary varieties and their projective
embeddings, descend to a field finitely generated over $k$. Spread
them over a smooth integral finite-type $k$-scheme $B$. Shrink $B$
so that the auxiliary families are smooth projective with
geometrically integral fibers, the cycles have their expected
codimensions, and the relative Chow equality holds. The last step
follows either by spreading its principal-divisor certificate or by
the direct-limit description of Chow groups over the function field.

Choose a closed point $b\in B$. Its residue field $E$ is finite over
$k$. Since $B$ is smooth, refined pullback to $b$ commutes with
the smooth pullbacks, proper projections and diagonal intersections
which define correspondence composition. The constant family
$X\times B$ specializes to $X_E$, while the auxiliary families
specialize to smooth curves and surfaces. Thus the diagonal
certificate specializes with all its degree shifts.
\end{proof}

\begin{theorem}[Complex certificates and surface moduli]\label{thm:moduli}
Let $X/F$ be smooth projective over a number field. If for one complex
embedding $h(X_\C)$ retracts through products of curves and surfaces
in rational Chow motives, then $X$ has a full height domain, also
after every finite number-field extension. In particular this holds
if $X_\C$ is a smooth projective moduli space appearing in
B\"ulles's theorem: Gieseker-stable twisted sheaves, or stable objects
for a generic stability condition, on a complex projective K3 or
abelian surface. Finite products of these spaces and the varieties
in \thmref{thm:hilbert} are allowed, with an everywhere potentially
good projective coefficient.
\end{theorem}

\begin{proof}
By \lemref{lem:specialize-certificate}, a finite extension $E/F$
carries an actual diagonal certificate of the required kind.
Complete it at a place above each finite place of $F$ and apply
\corref{cor:surface-envelope}; full-degree descent gives the
original local kernel vanishing. The same construction over a
compositum treats every new cycle after finite extension. Apply
\lemref{lem:two-degrees}.

For the moduli spaces, \cite[Theorem 0.1]{Bulles} supplies precisely
an actual rational Chow retraction through powers of the underlying
surface. The first assertion applies. For products, tensor the
certificates; the intermediate varieties remain products of curves
and surfaces. The surface's finite-dimensionality is not an
assumption of the cited retraction theorem.
\end{proof}
